\section{O formato de cada indicador}
\label{sec:distribuicoes}

Cada indicador tem uma forma própria de variar entre os hospitais, e
conhecer essa forma tem uma consequência prática para a gestão: por trás
de um histograma de aparência estranha costuma haver um fato concreto da
rede --- dois perfis de hospital convivendo lado a lado, um acúmulo de
zeros com significado, uma cauda que exige leitura em escala própria.
Esta seção reúne esses achados de formato. A justificativa estatística
formal --- qual família de distribuição descreve cada indicador e o
quanto ela se ajusta aos dados, com as tabelas de nota de ajuste e os
gráficos QQ correspondentes --- foi reunida no
Apêndice~\ref{ap:diagnosticos}, para não interromper a leitura.

\subsection{Mortalidade: uma proporção com dois grupos de hospitais}

A mortalidade é uma proporção (fica entre 0 e 1), e o molde clássico
para proporções é a distribuição \textbf{Beta}. O ajuste de uma Beta
única descreve bem a escala geral, mas erra o pico de hospitais com
mortalidade quase nula. A Figura~\ref{fig:mistura} mostra a solução e,
ao mesmo tempo, a explicação: somando \textbf{duas Betas} (uma para
cada grupo de hospitais), a curva acompanha o histograma com nitidez. O
ponto central para a gestão é que esses dois grupos
\textbf{existem de verdade no painel}: um grupo de \textbf{35 hospitais
de baixo óbito estrutural} (mortalidade mediana abaixo de 2\%; 16 são
hospitais especializados e os outros 19, hospitais gerais de baixo
volume) e os \textbf{254 demais}, hospitais gerais com mortalidade em
torno de 6\%.
A aparente anomalia do histograma é, portanto, apenas o retrato de dois
perfis assistenciais diferentes convivendo na rede. Como também existem
observações com mortalidade exatamente zero (0,9\%), o molde completo
para a mortalidade é a \textbf{Beta com um tratamento especial para os
zeros}.

\comoler{no painel da esquerda, compare as duas curvas com as barras: a
tracejada passa por cima do vale e erra os dois picos; a escura desenha
os dois picos. No painel da direita, as curvas tracejadas coloridas são
os dois grupos de hospitais (baixo óbito à esquerda, demais à direita)
e a curva preta é a soma dos dois, praticamente idêntica ao histograma:
a forma estranha da mortalidade é só a mistura de dois perfis de
hospital.}

\begin{figure}[H]
\centering
\subfig{fig_ae_12_mistura_mort_all.png}{Um molde só vs a soma de dois}
\subfig{fig_ae_12_decomposicao_mort_all.png}{Os dois grupos reais de
hospitais}
\caption{Mortalidade geral: à esquerda, a curva cinza tracejada é o
molde único (que erra o pico perto de zero) e a curva escura é a soma
de dois moldes, que acompanha o histograma; à direita, a mesma soma
reconstruída a partir dos dois grupos reais de hospitais do painel.}
\label{fig:mistura}
\end{figure}

\comofeito{os dois painéis têm no eixo horizontal a mortalidade geral
(proporção, entre 0 e 1) e no vertical a quantidade das 3.179
observações em cada faixa. À esquerda, sobre o histograma, foram desenhadas duas
curvas ajustadas: uma Beta única e a soma de duas Betas. À direita, a
soma de duas Betas foi reconstruída separando de fato os hospitais nos
dois grupos observados no painel (baixo óbito estrutural e demais), pela
mortalidade mediana de cada hospital.}

\subsection{Alta complexidade: o peso dos zeros}

A fração de alta complexidade também é uma proporção, mas com uma
particularidade dominante: \textbf{um terço das observações vale
exatamente zero} (hospitais que não realizam alta complexidade) e uma
pequena parcela vale exatamente um (hospitais que só realizam alta
complexidade). Esses dois acúmulos nas pontas exigem um tratamento
estatístico próprio --- o modelo conhecido pela sigla \textbf{ZOIB},
cuja construção e ajuste estão no Apêndice~\ref{ap:diagnosticos}. Para a
gestão, o que importa é que \textbf{os zeros não são aleatórios}
(Figura~\ref{fig:zeros}): a proporção de hospitais sem alta complexidade
cai de 80,1\% na faixa 2 de complexidade para zero na faixa 6, e de
37,5\% nos Filantrópicos para 10,7\% nas OSS. Ou seja, saber a
complexidade, o porte e a categoria do hospital já diz muito sobre a
chance de ele estar no grupo do zero.

\comoler{cada barra é a proporção de hospitais-ano daquele grupo que
não realizou nenhuma internação de alta complexidade. A queda das
barras da esquerda para a direita no painel (b) mostra que o zero é
questão de perfil e estrutura: quase todos os hospitais de faixa 2
estão no zero, e nenhum de faixa 6 está.}

\begin{figure}[H]
\centering
\subfig{fig_ae_15_zeros_pct_alta_categoria.png}{Por categoria de
gestão}
\subfig{fig_ae_15_zeros_pct_alta_faixa.png}{Por faixa de complexidade}
\caption{Percentual de observações com alta complexidade exatamente
zero, por grupo de hospitais.}
\label{fig:zeros}
\end{figure}

\comofeito{gráfico de barras. O eixo horizontal traz os grupos de
hospitais --- as categorias de gestão no painel (a) e as faixas
Barcelona de complexidade no painel (b) --- e o eixo vertical, o
percentual de observações daquele grupo em que a fração de alta
complexidade é exatamente zero. Cada barra é, portanto, a parcela de
hospitais-ano do grupo que não realizou nenhuma internação de alta
complexidade.}

\subsection{Faturamento e tempo de permanência: valores com cauda
longa}

Faturamento por saída e TMP são valores positivos com cauda à
direita\footnote{Guia dos tamanhos de amostra que
aparecem no documento: \textbf{3.179} é o painel completo;
\textbf{3.178} exclui a única observação com TMP igual a zero (Hospital
Infante D.\ Henrique, 2021), incompatível com a escala logarítmica;
\textbf{3.177} exclui as duas observações de
denominador frágil da Seção~\ref{sec:qualidade}, regra usada nos
modelos da fase de estimação; e \textbf{2.482} são as observações com
UTI ativa depois de retiradas essas duas frágeis. Cada conta foi
conferida no painel.}, e a leitura prática que interessa à gestão é uma
só: \textbf{os dois indicadores devem ser analisados em escala
logarítmica}. Nessa escala, o faturamento vira um morro quase simétrico
(Figura~\ref{fig:histlog}) e o TMP hoje se ajusta quase perfeitamente,
porque a consolidação do painel (ETAPA F) removeu os hospitais de longa
permanência que, na versão anterior, criavam um teto de 30,5 dias que
distorcia a distribuição. A comparação formal entre as famílias de
distribuição candidatas --- com a nota de ajuste de cada uma e os
gráficos correspondentes --- está reunida no
Apêndice~\ref{ap:diagnosticos} (Tabela~\ref{tab:familias}), e confirma
que faturamento e TMP são bem descritos pelo molde log logístico nessa
escala.

\comoler{na escala logarítmica, o faturamento (a) vira um morro quase
simétrico, sinal clássico de que a análise deve ser feita nessa escala.
No TMP (b), o morro principal fica entre 3 e 7 dias, com uma cauda
curta à direita e sem o segundo grupo isolado de longa permanência que
o painel anterior exibia.}

\begin{figure}[H]
\centering
\subfig{fig_ae_03_hist_log_custo_saida.png}{Faturamento por saída}
\subfig{fig_ae_03_hist_log_tmp.png}{TMP}
\caption{Distribuições de faturamento e TMP com o eixo horizontal em escala
logarítmica (os rótulos mostram os valores reais).}
\label{fig:histlog}
\end{figure}

\comofeito{são os mesmos histogramas de faturamento por saída e TMP, mas
com o eixo horizontal em \textbf{escala logarítmica} (as marcas avançam
por multiplicação, e os rótulos mostram os valores reais em reais e em
dias); o eixo vertical segue acompanhando a quantidade de observações
por faixa. A escala
log foi escolhida justamente porque esses dois indicadores têm cauda
longa à direita, que na escala linear comprime a maioria dos hospitais
num só bloco.}

\subsection{Produção: contagens muito mais variáveis que o usual}

A produção anual (número de saídas) é uma contagem. O molde mais
simples para contagens, a Poisson, supõe que a variação entre hospitais
seja comparável à média; aqui \textbf{a variação é milhares de vezes
maior que a média} (razão de 6.249 no agregado e de 102 mesmo comparando
cada hospital só consigo mesmo ao longo dos anos). O molde adequado é a
\textbf{Binomial Negativa}, que aceita essa variação extra.
